\section{Problem Statement}

\subsection{Introduction}
In the context of the increasing prevalence of speech recognition technology, processing audio signals in real-world environments remains a significant challenge due to the impact of unwanted noise. This research focuses on applying Canonical Correlation Analysis as a powerful preprocessing technique to enhance the input quality for Speech-to-Text systems. Unlike traditional noise filtering methods that often focus on frequency suppression, CCA aims to identify and establish links between multi-dimensional data sources. This allows the system to recognize and preserve the core invariant features of human speech, even when they are obscured by complex noise components.

The primary objective of this experiment is to use CCA to learn these invariant features, which are the most stable and essential signal components of speech. By maximizing the correlation between reference datasets, the study aims to build an intelligent filter capable of automatically separating the target signal from random environmental noise without requiring strict assumptions about the type of noise involved. The final result is a preprocessing solution that significantly improves the accuracy of the output text, ensuring the robustness of the recognition system even when the input audio quality is severely degraded by external factors.

\subsection{Input}
The experimental procedure was designed in two distinct phases to optimize recognition capabilities.

\subsubsection{Training Phase}
The experimental workflow begins with a training phase designed to build an optimal signal separation model using paired reference data. During this stage, the system simultaneously processes two input datasets:
\begin{itemize}
    \item \textbf{Clean Dataset ($X^{(1)}$):} Consists of standard audio samples recorded in an ideal, noise-free environment.
    \item \textbf{Noisy Dataset ($X^{(2)}$):} Generated by integrating Additive White Gaussian Noise (AWGN) directly into the clean samples.
\end{itemize}

Simulating noise in this manner allows for precise control over data variability, creating a complete experimental environment for the algorithm to learn the mapping between the clean entity and its noisy counterpart. The core role of using the clean and noisy data pair is to provide the CCA method with a precise frame of reference to define the stable mathematical structure of human speech. Through the mechanism of maximizing covariance, CCA learns to focus on the components with the highest correlation between the two datasets (the target signal) while automatically identifying and suppressing uncorrelated random fluctuations (the noise components). The outcome of this phase is an optimized set of canonical weights, acting as a feature map that enables the system to retrieve useful information from chaotic data sources in subsequent execution stages.

\subsubsection{Testing Phase}
Once the optimal canonical weights have been established during the training phase, the system is deployed for real-world recognition tasks using individual audio sources. In this phase, the input data can be any environmental audio signal, which is often degraded by ambient noise or white noise. 

Instead of attempting to reconstruct a clean acoustic waveform, the system applies the learned projection matrices to transform the noisy audio data into a new feature space. Within this space, CCA functions by filtering out irrelevant noise fluctuations and preserving only the signal components that correlate most strongly with the learned linguistic structures. This process cleans the data at the feature level, creating optimal conditions for subsequent Deep Learning models to analyze and extract content with maximum precision.

\subsection{Output}
This approach delivers highly reliable recognized text with strong fidelity to the original speech content, even in noisy and challenging environments. 

The final result is not a denoised audio waveform, but a robust speech representation where essential linguistic information is preserved and noise influence is minimized. This refined representation allows the decoding stage to produce accurate transcription decisions, ensuring the system can still recover the correct utterance in scenarios where conventional methods would fail.